% ============================================================
%  CHƯƠNG 2: CƠ SỞ LÝ THUYẾT
% ============================================================

\section{[Tên lý thuyết / khái niệm cơ bản thứ nhất]}

[Trình bày các khái niệm, lý thuyết nền tảng cần thiết để hiểu phương pháp
đề xuất trong luận văn. Nên có trích dẫn đầy đủ cho mỗi khái niệm.]

\subsection{[Tiểu mục]}

[Nội dung chi tiết.]

% -------------------------------------------
\section{[Tên lý thuyết / công cụ thứ hai]}

\subsection{[Tiểu mục]}

[Nội dung chi tiết. Ví dụ về bullet list:]

Các đặc điểm chính bao gồm:
\begin{itemize}
  \item Đặc điểm thứ nhất: [mô tả]
  \item Đặc điểm thứ hai: [mô tả]
  \item Đặc điểm thứ ba: [mô tả]
\end{itemize}

% -------------------------------------------
\section{[Các phương pháp / mô hình liên quan]}

\subsection{[Phương pháp A]}

[Mô tả phương pháp A, cách hoạt động, ưu nhược điểm.]

\subsection{[Phương pháp B]}

[Mô tả phương pháp B.]

% Ví dụ công thức
Hàm mất mát cross-entropy được định nghĩa như sau:

\begin{equation}
  \mathcal{L}(\theta) = -\frac{1}{N} \sum_{i=1}^{N}
    \left[ y_i \log \hat{p}_i + (1-y_i) \log(1-\hat{p}_i) \right]
  \label{eq:cross_entropy}
\end{equation}

Trong đó: $y_i \in \{0,1\}$ là nhãn thực tế; $\hat{p}_i$ là xác suất dự đoán
của mô hình; $N$ là tổng số mẫu huấn luyện.

% -------------------------------------------
\section{Tóm tắt chương}

[Tóm tắt ngắn gọn những nội dung lý thuyết đã trình bày trong chương 2
và mối liên hệ với phương pháp đề xuất ở Chương 3.]
